\section{Materials and Methods}

\subsection{Experimental Setup}
\begin{figure}[ht]
\includegraphics[width=\textwidth]{Fig1.png}
\caption{\textbf{Experimental setup and protocol.} \textbf{A}, Study 1 setup with participant wearing headphones (blue cord) and chest-mounted vibrotactile stimulator (red cord) with eye shade for visual deprivation. \textbf{B}, Study 2 setup with breath-responsive VR headset providing visual feedback of respiratory state. \textbf{C}, Box breathing protocol (4s inhale, 4s hold, 4s exhale, 4s hold) with looming sounds and tactile stimuli presented during breath-hold phases. D1-D5 represent stimulus onset asynchronies between audio and tactile stimuli.}
\label{study5:fig:exp_setup}
\end{figure}

Participants will be seated comfortably while wearing headphones and a chest-mounted vibrotactile stimulator (Fig. \ref{study5:fig:exp_setup}). For Study 1, participants will wear an eye shade to eliminate visual input (Fig. \ref{study5:fig:exp_setup}A), while in Study 2, participants will wear a VR headset that provides exteroceptive feedback about their respiratory state (Fig. \ref{study5:fig:exp_setup}B). In both studies, participants will follow a guided box breathing protocol: inhalation (4s), inhalation hold (4s), exhalation (4s), and exhalation hold (4s). During the breath-hold phases, tactile stimuli will be delivered along with looming sounds at various stimulus onset asynchronies (D1-D5), as shown in Fig. \ref{study5:fig:exp_setup}C. Looming sounds will be presented via headphones, simulating objects approaching from either side. The vibrotactile stimulator on the chest will deliver tactile stimuli synchronized with the sounds at various temporal delays. Participants will press a button when they detect the vibrotactile stimulus. Respiratory compliance will be monitored using a respiratory belt transducer placed around the lower ribcage.

\subsection{Study Design}
The experimental paradigm employs a within-subjects design investigating peripersonal space boundaries between respiratory phases, most closely following established protocols used to test PPS modulation in response to vestibular stimulation \citep{pfeiffer2018vestibular}. To minimize time constraints while ensuring participants maintain box breathing throughout the experimental protocol, we implemented 5 stimulus onset asynchronies (SOAs) \citep{lee2021altered, mine2022adaptation, teneggi2013social} with 12 repetitions per SOA-condition combination \citep{noel2015peripersonal, galli2015wheelchair}. While other PPS studies have employed more extensive designs with 7 SOAs and 21 repetitions, our 8-second trial duration necessitated by the box breathing protocol created practical time constraints that required using these parameters to maintain a reasonable total experiment length.

Participants follow continuous breathing instructions: inhale for 4 seconds, hold breath for 4 seconds, exhale for 4 seconds, hold breath for 4 seconds. During this cycle, they receive auditory and/or tactile stimuli to which they respond via mouse press. Participants are instructed to respond only to vibrotactile stimuli while ignoring looming sounds. The experiment consists of three primary conditions (each comprising approximately one-third of all trials): Inhale (tactile stimulus during inhalation hold with accompanying auditory stimulus), Exhale (identical protocol during exhalation hold), and Baseline (tactile stimulus without auditory stimulus, evenly distributed between respiratory phases). Additionally, 10\% of trials are designated as catch trials, featuring auditory stimuli without tactile stimulation, requiring response inhibition.

\begin{figure}[ht]
\includegraphics[width=\textwidth]{trial_design.png}
\caption{\textbf{Experimental design showing trial distribution.} This figure shows trial distribution across conditions and distances. The Inhale and Exhale conditions each comprise 12 repetitions per SOA, while Baseline includes 6 repetitions per respiratory phase. Catch trials feature 2 repetitions per respiratory phase per SOA. Total: 200 trials (26.7 minutes).}
\label{study5:fig:trial_design}
\end{figure}

To ensure attentional vigilance and prevent response habituation, we incorporated 20 catch trials evenly distributed across all SOAs, with 2 trials per respiratory phase per SOA. This design results in a total of 200 trials (60 trials each for Inhale, Exhale, and Baseline conditions, plus 20 catch trials). Each trial lasts 8 seconds, resulting in a total experimental duration of approximately 26.7 minutes excluding breaks and instructions. Each participant received a unique randomized trial sequence with constraints preventing more than three consecutive trials of the same type. We verified randomization quality through chi-square tests confirming balanced distribution across breath phases and analysis of consecutive trial patterns. Additional temporal jitter in stimulus timing further enhanced unpredictability of the stimulus onset in each trial.

\subsection{Auditory-Tactile Stimulus Presentation}
To generate the auditory looming stimuli, we created binaural sound files using four types of noise---pink, white, brown, and blue---selected to counteract sound habituation, maintain stimulus novelty, and avoid repetition inhibition effects throughout the experiment. All stimuli were generated at a 44.1 kHz sampling rate with a 3-second duration. The looming trajectory was simulated as a frontal approach ($0^\circ$ azimuth, $0^\circ$ elevation) with the sound source beginning at 2.0 meters distance and approaching to 0.5 meters, maintaining a constant velocity throughout the motion path. We applied Head-Related Transfer Functions (HRTFs) from the FABIAN head model (neutral head orientation) \citep{brinkmann2017fabian} to simulate spatial filtering by the head and pinnae across different positions. This implementation utilized the SOFA (Spatially Oriented Format for Acoustics) file format via the sofar library. We applied an inverse square law (attenuation factor = 40) to modulate stimulus amplitude while calculating precise source-to-ear distances based on an 18 cm head diameter. Each noise type was normalized to 80 percent of maximum amplitude to prevent digital clipping while maintaining consistent loudness across stimulus types. The final stimuli incorporated a gentle fade-out (100 ms) at the end to avoid auditory artifacts.

For tactile stimulation, we used a Woojer haptic belt to deliver stimuli to participants' lower chest while measuring peripersonal space boundaries. The chest strap acts as an audiotactile transducer, and we used 200 Hz sinusoidal wave (amplitude +4dB) combined with a 50 Hz decay signal (amplitude -22dB) for 100ms to deliver brief vibrations. In order to deliver auditory and tactile stimuli in a synchronized manner, for each participant, we combined the randomized sequence of auditory looming stimuli and tactile stimuli into a single stereo WAV file overlaid onto the breathing instructions in the background, with looming sounds in the left channel and tactile trigger signals in the right channel. This combined file was then played through a Komplete Audio interface, which routed the left channel to Output 1/2 (connected to headphones for auditory stimulation) and the right channel to Output 3/4 (connected to the haptic belt for tactile stimulation). This approach ensured perfect temporal synchronization between auditory and tactile stimuli, as both signals originated from the same audio file and were processed through the same device. The ASIO driver provided low-latency performance essential for maintaining precise timing relationships between the multisensory stimuli.

\subsection{Participants}
Thirty healthy adults (18-40 years) will participate after screening for respiratory and neurological disorders.
